\scanpage{21}
\mspage{- IV-\handwritten{\unsure{952|953}} -}
\uline{neau} $({}^{\mathrm h}A)_{\underline p}/\underline p.({}^{\mathrm h}A)_{\underline p}$ \uline{est} \deleted{\illegible{1}} \uline{composé direct d'un nombre fini d'extensions
algébriques séparables de} $\handwritten{\kappa}(\underline p)$ (ce qui entraîne que les fibres du morphisme
canonique $\operatorname{Spec}({}^{\mathrm h}A)\to\operatorname{Spec}(A)$ sont \uline{géométriquement régulières}) .

(i) Comme $ {}^{\mathrm h}A$ est un $A$-module fidèlement plat,il résulte de (2.1.\unsure{13}) que si
$ {}^{\mathrm h}A$ est \deleted{\illegible{1}} réduit (resp. normal) , il en est de même de $A$ . Pour prouver la
réciproque , on peut se borner au cas où l'anneau $A$ est local , de sorte \deleted{\illegible{1}}
qu'avec les notations de (18.6.4) , on a $ {}^{\mathrm h}A=\varinjlim A_{\handwritten{\alpha}}$ . Or , il résulte de la
définition des $A_{\handwritten{\alpha}}$ et de (17.5.7) que si $A$ est réduit (resp. intègre et intégralement
clos) , il en est de même des $A_{\handwritten{\alpha}}$ ; \deleted{\unsure{xx}} en outre , les \deleted{\illegible{1}
\illegible{5}} morphismes \deleted{$\operatorname{Spec}(A_{\handwritten{\beta}})\to\operatorname{Spec}(A_{\handwritten{\alpha}})$} \inserted{\handwritten{$A_\alpha\unsure{-|\to}A_\beta$}} pour $\handwritten{\alpha\leq\beta}$
% slanted
\margin{\handwritten{\unsure{pinceaux|précisions} \unsure{mélangés|mélangées}!}}
sont \deleted{surjectifs} \inserted{\handwritten{injectifs}} en vertu de \deleted{(2.5.4,(ii))}\inserted{\handwritten{\unsure{----}}} . On conclut donc de (5.13.2 ) (resp
(5.13.4 )) que $ {}^{\mathrm h}A$ est réduit (resp. intègre et intégralement clos) .

(ii) \inserted{On peut se borner au cas où $A$ est local .} En vertu du fait que le foncteur $\varinjlim$ permute au produit tensoriel ,
la fibre du morphisme $\operatorname{Spec}({}^{\mathrm h}A)\to\operatorname{Spec}(A)$ en un point $p$ est limite inductive
des fibres des morphismes $\operatorname{Spec}(A_{\handwritten{\alpha}})\to\operatorname{Spec}(A)$ en ce point . . Comme $ {}^{\mathrm h}A$ est
noethérien , on voit , compte tenu de (17.6.2) , que l'on est ramené à prouver
le lemme suivant :

\uline{Lemme} (18.6.9.1) .- \uline{Soient} \deleted{\illegible{1}} $(B_{\handwritten{\alpha}})$ \uline{un système inductif d'anneaux artiniens} ,
$B=\varinjlim B_{\handwritten{\alpha}}$ \uline{sa limite inductive . Si} $B$ \uline{est noethérien , alors} $B$ \uline{est artinien} ;
\deleted{\illegible{6}} \uline{si de plus , pour} $\handwritten{\alpha\leq\beta}$ , \uline{les
homomorphismes} \inserted{\handwritten{$\varphi_{\beta\alpha}:$}} $B_{\handwritten{\alpha}}\to B_{\handwritten{\beta}}$ \uline{sont injectifs , alors il existe} $\handwritten{\lambda}$ \uline{tel que , pour} $\handwritten{\alpha\geq\lambda}$
\uline{le nombre des composants locaux} \deleted{\illegible{1}} $B_{\handwritten{\alpha}}^{(i)}$ \uline{de} $B_{\handwritten{\alpha}}$ \uline{soit constant} , \deleted{\illegible{1}}
\uline{que} $\handwritten{\varphi_{\beta\alpha}}(B_{\handwritten{\alpha}}^{(i)})\subset B_{\handwritten{\beta}}^{(i)}$ \uline{pour tout $i$ et que les} $B^{(i)}=\varinjlim B_{\handwritten{\alpha}}^{(i)}$ \uline{soient les composant\handwritten{s}
locaux de} $B$ .

Soient en effet $Y_{\handwritten{\alpha}}=\operatorname{Spec}(B_{\handwritten{\alpha}})$ , $Y=\operatorname{Spec}(B)$ , qui , en tant qu'espace topologique
, est égal à $\varprojlim Y_{\handwritten{\alpha}}$ (8.2.10) . Les espaces $Y$ étant \deleted{\illegible{1}} \inserted{\handwritten{finis et}} discrets , la seconde
assertion du lemme résultera de la première et du fait que si les $\handwritten{\varphi_{\beta\alpha}}$
% slanted
\margin{\handwritten{\unsure{pas très clair}}}
sont injectifs , les morphismes canoniques $Y_{\handwritten{\beta}}\to Y_{\handwritten{\alpha}}$ sont surjectifs . On est
ainsi ramené à prouver le lemme purement topologique suivant :

\uline{Lemme} (18.6.9.2) .-\uline{Soit} $(Y_{\handwritten{\alpha}})$ \uline{un système projectif d'espaces topologiques finis
et discrets . Si} $Y=\varprojlim Y_{\handwritten{\alpha}}$ \uline{est noethérien} , $Y$ \uline{est fini et discret} .
% archive sheet number: 20
